\documentclass[12pt]{article}
\usepackage[utf8]{inputenc}
\usepackage[T1]{fontenc}
\usepackage{hyperref}
\title{Masela, East Language Resources for K-12 Curriculum Development}
\author{Compiled by Claude Sonnet 5 (claude-sonnet-5) for Andr\'as Kornai}
\date{}
\begin{document}
\maketitle
\section{Introduction}
East Masela (ISO 639-3 vme, Glottocode east2462) is an Austronesian language of Marsela Island in the Babar Islands, southern Maluku, Indonesia. Speaker estimates vary (about 520 per Verbix, about 1,000 per Joshua Project).
\section{Primary Educational Resources}
None found.
\section{Linguistic and Reference Materials}
Masela language overview: \url{https://en.wikipedia.org/wiki/Masela_language}. Glottolog: \url{https://glottolog.org/resource/languoid/id/east2462.bigmap.html}. Verbix page: \url{https://docs.verbix.com/Languages/MaselaEast}. No grammar or dictionary specific to East Masela was found.
\section{Religious and Community Resources}
Joshua Project people-group profile: \url{https://joshuaproject.net/people_groups/13567}.
\section{Digital and Audio Resources}
None found.
\section{Curriculum Development Context}
Documentation appears minimal. Any curriculum would require fieldwork and community partnership, with Indonesian as the school language.
\end{document}